\clearpage\thispagestyle{empty}
\pdfbookmark[0]{Quellen zum Ceres-Bericht vom April 1802}{ceres180204-sources}
\begin{center}\textbf{Quellen zum Ceres-Bericht vom April 1802}\end{center}
\noindent Diese KI-Übertragung in heutiges Deutsch folgt dem deutschen Auszug in Carl Friedrich Gauß’ \href{https://gdz.sub.uni-goettingen.de/id/PPN236007467}{\textit{Werken}, Band~VI} (Göttingen, 1874), S.~209--210. Franz Xaver von Zach berichtet über Ceres und gibt Rechnungen von Carl Friedrich Gauß sowie Durchmesserbeobachtungen von Johann Hieronymus Schröter wieder.
\par\smallskip\noindent Von Zachs \href{https://api.digitale-sammlungen.de/iiif/image/v2/bsb10538597_00391/full/1450,/0/default.jpg}{ursprünglicher Aufsatz}, Nr.~XL, erschien im Aprilheft 1802 der Zeitschrift \textit{Monatliche Correspondenz zur Beförderung der Erd- und Himmels-Kunde} (Gotha, Becker), Band~V, S.~379--404. Der Auszug in den Werken beruht auf \href{https://api.digitale-sammlungen.de/iiif/image/v2/bsb10538597_00401/full/1450,/0/default.jpg}{S.~389--390} und \href{https://api.digitale-sammlungen.de/iiif/image/v2/bsb10538597_00414/full/1450,/0/default.jpg}{S.~402--404}. Auf der \href{https://api.digitale-sammlungen.de/iiif/image/v2/bsb10538597_00416/full/1450,/0/default.jpg}{letzten Seite} steht auch der Schluss des ursprünglichen Aufsatzes. Die Digitalisate der Zeitschrift stammen von der Bayerischen Staatsbibliothek.
\par\smallskip\noindent Übertragung und Satz wurden mit KI erstellt.
\clearpage
